\section{Introduction}
In many modern applications, the quantity of scientific interest is expensive to measure, while a cheap surrogate is available for nearly every unit. Cheapness does not settle the inferential question, because a surrogate that predicts well on average can be inaccurate in the neighbourhood or in the tail where inference is required. In the air-monitoring study below, a linear calibration trained on historical collocations leaves a mean difference of $-0.191$ micrograms per cubic metre near 30\% relative humidity in the subsequent audit period, while the calibrated sensor assigns 6.068\% of local observations above a concentration of 15 against 9.307\% for the reference. Mean calibration therefore leaves the distributional question open: which of these discrepancies can be established from a limited reference sample.

We observe covariates $X$ and a surrogate $S$ for most units, a reference outcome $Y$ only in a verified subset, and we target the conditional distribution of $Y$ at a chosen $x_0$. Augmented inverse-probability weighting remains valid when the fitted outcome regression is inaccurate, and it joins a surrogate-free anchor to a surrogate-assisted estimator through a common borrowing path. The question we ask is which feature of uncertainty should determine the weight on that path.

For a simultaneous band, average variance is the wrong criterion. The band's critical value depends on the joint covariance across thresholds, including how the correction redistributes variance. Proposition~\ref{prop:separation} gives a valid two-threshold path whose population trace-minimising weight reduces total variance and increases the simultaneous critical value, so an exactly known optimal average-variance weight can lengthen the band. Estimating a band-optimal weight adds a second difficulty, because a small apparent reduction may be sampling noise in the covariance estimate.

The observed-data representation is classical. Surrogate-outcome and two-phase methods establish it \citep{Pepe1992,RobinsRotnitzkyZhao1994,TaoZengLin2017,Tsiatis2006}, rejective two-phase designs quantify its efficiency gains \citep{YangDing2025}, and combining probability with non-probability samples extends it to modern data sources \citep{YangKimSong2021}. Semi-supervised inference supplies the complementary lesson that unlabelled units can sharpen inference for prediction performance and regression parameters \citep{GronsbellCai2018,GronsbellEtAl2022,AzrielEtAl2022,ChakraborttyCai2018,ZhangBrownCai2019,Song2024}, and modern learning with missing covariates shows when such gains persist under flexible estimators \citep{MaWangSamworth2026}. Adaptive augmentation optimises auxiliary-information use for parameter estimation \citep{CaoTsiatisDavidian2009}. These procedures choose auxiliary information for parameters, and the selection criterion for a threshold-indexed local distribution has remained open.

Prediction-powered inference has made the label-scarce setting a central object of study. Its basic construction and efficient variants \citep{AngelopoulosEtAl2023,AngelopoulosDuchiZrnic2023,ZrnicCandes2024} treat a fitted predictor as an auxiliary variable, and follow-up work examines label contamination, surrogate-powered regularisation, adaptive shrinkage and post-prediction inference \citep{SesiaEtAl2025,JiLeiZrnic2025,ChenEtAl2025,LiIgnatiadis2025,MiaoEtAl2025}. The finite-reference cost of estimating borrowing weights is now documented \citep{mani2025nofreelunch,Eyre2025}, and \citet{KallusMao2025} characterise the efficiency gain from surrogates in estimation with limited outcome data. These results certify or regularise a fitted predictor for scalar targets; a weight chosen for a simultaneous band on a whole conditional distribution requires a different criterion.

Distributional targets have their own line of work. \citet{SuiEtAl2026} study prediction-powered conditional functionals, \citet{TestaEtAl2025} allow selective labels with decreasing overlap, and \citet{Wang2026} develops surrogate-assisted distributional inference for transported outcomes. Local distribution estimation with fully observed outcomes is established \citep{CattaneoJanssonMa2024,CattaneoEtAl2024}, conformal methods deliver distribution-free prediction sets for counterfactuals, individual effects and survival times \citep{LeiCandes2021,CandesLeiRen2023}, and generation-powered and multi-source extensions develop functional borrowing and confidence-region volume \citep{zhang2026generation,li2026multisource}. Rearrangement handles crossing quantiles \citep{ChernozhukovEtAl2010}. Anderson's covariance-dominating operator improves Gaussian probabilities of centred convex symmetric sets \citep{Anderson1955}, and a fitted path need not admit such an ordering.

Guarantee calibration, selection control and transfer complete the picture. Local weights and randomization trade bias for localised validity in conformal prediction \citep{HoreBarber2025}, conditional and selection-conditional coverage are characterised with matching impossibility results \citep{GibbsEtAl2025,JinRen2025}, and prediction sets adapt to unknown covariate shift through doubly robust calibration \citep{QiuEtAl2023,YangEtAl2024}. Error-controlled selection with subsampling \citep{ShahSamworth2013} and covariate-powered testing \citep{IgnatiadisHuber2021} control the selection error itself, while doubly robust testing with generative models inserts learned components into test statistics \citep{ZhangEtAl2026}. Transfer theory quantifies when auxiliary data help and when they harm, in high-dimensional and nonparametric regression and classification \citep{LiCaiLi2022,GuHanDuan2025,AuddyEtAl2026}. Our negative-transfer question is local: a surrogate correction can help the anchor at some thresholds and hurt it at others.

We evaluate the resulting procedure in two settings. A simulation study compares six borrowing criteria under alignment, weak surrogates, local reversal, absent signal and an exact zero path, then isolates the contribution of pairing in a matched ablation, tests dense threshold grids and non-Gaussian outcomes, and evaluates exact operating characteristics in a binary local experiment. An air-monitoring application separates historical mean calibration from subsequent distributional verification: all 2019 site-days fit and freeze the calibration and the candidate regressions, and designed Bernoulli verification on the 2020--2021 audit frame measures which humidity-specific exceedance discrepancies a given reference budget can establish.
We develop \emph{paired gain certification} for a local distribution. We estimate the difference between the anchor and candidate simultaneous critical values under a common covariance, and its uncertainty vanishes at zero borrowing. We exploit that cancellation to obtain an oracle bound proportional to the beneficial borrowing weight, an implementable multiplier calibration with independent numerical guards, and simultaneous inference after selection without a unique optimal weight. A binary local experiment gives the reference-information scale for certifying weak improvements. The general learning-and-testing principle is established \citep{angelopoulos2021learn}; the contribution here is the paired critical-value geometry, its local information scaling, and its use in distributional inference.

Section~2 defines the estimator and algorithm, Section~3 states the certification, inference and information results, Section~4 compares borrowing criteria and evaluates their extensions, and Section~5 applies the method to dated air-monitoring collocations before the concluding discussion.

